\section{Cyclic oxides: steam compatibility and recovery of selective oxygen delivery}
\label{sec:cyclic}

\textbf{Proposed contribution.} Test whether replacing the post-oxidation Ar purge with a selected steam-rich purge changes useful function in Li\textsubscript{2}CO\textsubscript{3}-coated La\textsubscript{0.8}Sr\textsubscript{0.2}FeO\textsubscript{3} (LSF) under ordinary cycling without supplementary CO\textsubscript{2}, then test recoverability and whether concurrent CO\textsubscript{2} is more useful than delayed treatment. The prospective result is a compatibility and recovery comparison under specified histories, with a conditional prediction of complete-cycle ethylene output. Steam damage, lithium escape, and the benefit of protective CO\textsubscript{2} remain hypotheses. The post-ethane purge remains dry Ar and is outside the substitution claim.

\subsection{Why the question matters, and what prior work already establishes}

Ethane oxidative dehydrogenation produces ethylene while removing hydrogen as water. In chemical looping, an oxygen carrier supplies oxygen during the ethane step and is subsequently reoxidized, separating ethane conversion from the gaseous oxidant feed. The ideal selective reduction is
\begin{equation}
 \mathrm{C_2H_6+[O]_{carrier}\longrightarrow C_2H_4+H_2O}.
 \label{eq:cyclic-odh}
\end{equation}
The challenge is delivering enough oxygen for useful conversion without oxidizing hydrocarbons to CO or CO\textsubscript{2}. Gao and colleagues established a carbonate-coated LSF carrier and attributed its selectivity to a molten layer that blocks nonselective sites and transports peroxide species associated with Fe\textsuperscript{4+}/Fe\textsuperscript{3+} redox. They report approximately 0.42\,wt\% oxygen consumption over the first six minutes at 700\,$^\circ$C. The supported coating is reported to melt below bulk Li\textsubscript{2}CO\textsubscript{3}; hydroxide formation need not be the event that first makes it liquid.\cite{gao2020}

Brody and colleagues subsequently cycled 300\,g of coated LSF for more than 1200 elapsed hours. During the first 1000\,h at 700\,$^\circ$C, each cycle contained one minute of ethane feed, five minutes Ar purge, five minutes oxidation, and five minutes Ar purge. Average C\textsubscript{2+} carbon yield declined from 53.2\% over hours 0--300 to 47.24\% over hours 600--1000. C\textsubscript{2+} means ethylene plus C\textsubscript{3+} products, not ethylene alone. Later temperature and flow changes improved performance, so the later optimum was not sustained throughout the whole test. The methods specify 10\,wt\% carbonate whereas the conclusion says 20\,wt\%; actual loading must be measured.\cite{brody2022}

The process model in that study replaces the experimental Ar purges with steam heat-exchange steps. It does not measure how steam changes the coating.\cite{brody2022} This gap is the strongest reason to investigate steam compatibility. Steam cannot explain the decline already observed under Ar; explaining that decline is a separate aging problem. Nor does the source establish a continuously producing, thermally self-sufficient plant simply by observing exothermic half-cycles.

The chemical premise and much of the broader regeneration idea are also established. Barckholtz and colleagues observed carbonate/hydroxide responses to steam and CO\textsubscript{2} in a mixed Li/Na melt, although their Raman signals did not give absolute species concentrations.\cite{barckholtz2021} Chacko's dissertation reports actual mixed-carbonate/LSF cycling at 800\,$^\circ$C with 45\,wt\% salt and CO\textsubscript{2}/O\textsubscript{2} regeneration. It includes carbonate decomposition during inert purges and H\textsubscript{2} evolution during regeneration, attributed to reduced iron reacting with water released during hydroxide carbonation. These are close precedents for coupled coating/carrier regeneration, despite the different formulation and temperature.\cite[ch.~4.1]{chacko2025thesis}

Even prevention versus incomplete delayed recovery is prior art. Fereres and colleagues changed a bulk Li/Na/K carbonate melt from air to CO\textsubscript{2} after 30 days at 640\,$^\circ$C: further mass loss stopped without restoring the initial state under their treatment.\cite{fereres2018} The proposed originality is therefore specific: test the post-oxidation steam substitution and its recoverability in this lower-loading coating, assess whether a measured state predicts a new operating history, and determine whether an intervention repays its gas and time burden. Another demonstration that CO\textsubscript{2} stabilizes carbonate would be insufficient.

\subsection{The chemical hypothesis and its competing explanations}

A candidate component equilibrium is
\begin{equation}
 \mathrm{Li_2CO_3+H_2O\rightleftharpoons 2LiOH+CO_2},\qquad
 \frac{a_{\mathrm{LiOH}}^2}{a_{\mathrm{Li_2CO_3}}}
 =K(T)\frac{f_{\mathrm{H_2O}}}{f_{\mathrm{CO_2}}}.
 \label{eq:cyclic-hydrolysis}
\end{equation}
Here $a$ denotes condensed-phase activity and $f$ gas fugacity; the fugacity ratio approximates the partial-pressure ratio in the proposed atmospheric-pressure diagnostic. Lower steam/CO\textsubscript{2} ratio favors carbonate relative to hydroxide at fixed temperature. This direction does not supply a composition, equilibration time, or lithium-loss rate for supported LSF. Nonideal melt chemistry, substrate interactions, and changing coverage preclude importing a fuel-cell equilibrium constant uncritically. Barckholtz's experiments used hours per condition; they do not demonstrate a five-minute coating response.\cite{barckholtz2021}

Several outcomes could follow. A retained salt might change reversibly between chemical forms; it might redistribute away from important interfaces; lithium might move onto reactor contact materials or leave in transported material; or the carrier could change redox state or phase. These effects can occur sequentially. Lithium found downstream does not identify LiOH vapor, and unchanged total lithium does not establish unchanged coverage. A carbonate surface-signal decrease is not a quantitative lithium-loss measurement.

Restoration is also chemically coupled. The balanced reverse reaction,
\begin{equation}
 \mathrm{2LiOH+CO_2\longrightarrow Li_2CO_3+H_2O},
 \label{eq:cyclic-carbonation}
\end{equation}
releases water even with a dry inlet. Chacko's printed analogous reaction omits the factor of two before hydroxide; Eq.~\ref{eq:cyclic-carbonation} corrects that stoichiometry. The accompanying water can participate in carrier oxidation.\cite[ch.~4.1]{chacko2025thesis} Thus equal fed CO\textsubscript{2} or equal mean Fe valence cannot establish equivalent coating--carrier interfaces. The timing comparison measures a total treatment-order effect and restricts its interpretation using inventories and intermediate states.

\subsection{Experiment 1: ordinary-cycle compatibility, then a timing comparison}

Table~\ref{tab:cyclic-sequences} separates the operating comparison from the diagnostic and policy sequences. Ordinary oxidation remains part of every cycle; a special reset is an additional treatment.

\begin{table}[htbp]
\centering
\small
\caption{Sequence map for the proposed cyclic-oxide comparisons.}
\label{tab:cyclic-sequences}
\begin{tabularx}{\linewidth}{@{}>{\raggedright\arraybackslash}p{0.14\linewidth}>{\raggedright\arraybackslash}p{0.25\linewidth}>{\raggedright\arraybackslash}p{0.25\linewidth}>{\raggedright\arraybackslash}X@{}}
\toprule
Sequence & Exposure & Special-reset timing & Primary endpoint\\
\midrule
Ordinary operation & Five-minute post-oxidation steam-rich or Ar purge; fixed cycle block & None during the block; one terminal reset after operating output is measured & Pre-reset cycle and cumulative ethylene; separate post-reset recovery assay\\
Five-arm timing diagnostic & Two five-minute segments: wet/dry early/late CO\textsubscript{2}, plus Ar reference & One common reset after exposure, or after the fixed accumulation block & Post-reset ethylene and timing contrasts; intermediate state samples support attribution\\
Usable policy blocks & Concurrent protection, delayed restoration, untreated steam and Ar at the selected duty & Managed policies: frozen treatment cadence. Untreated/Ar: none. Terminal analytical reset follows counted output & Ethylene per complete time and resource demand; terminal function reported separately\\
\bottomrule
\end{tabularx}
\end{table}

Prepare one reproducible coating composition, quantify its initial Li content and phases, and establish stable short-cycle function with a dry Ar protocol. Use the methods-based nominal 10\,wt\% loading only as the starting recipe, documenting the measured final composition and loading convention. Begin the intervention immediately after ordinary oxidative regeneration. This selects the relatively oxidized post-oxidation purge; it does not yet represent the reducing post-ethane purge.

First compare ordinary cycling with a steam-rich post-oxidation purge against ordinary cycling with Ar, without supplementary CO\textsubscript{2} in either arm. Use 700\,$^\circ$C and atmospheric pressure, with the source's one-minute ethane, five-minute Ar, five-minute oxidation, and five-minute final purge as the starting schedule. Select one steam fraction in Ar, flow, and temperature trajectory from equipment qualification and the intended heat-exchange duty; hold the other cycle conditions fixed. Freeze that composition, the initial conditioning, cycle count, and functional averaging windows before comparing independent charges randomized within preparation and run-order blocks. Choose cycle and replicate counts from pilot variability, method uncertainty, and a predeclared consequential output loss, not by cycling until failure appears. Ordinary oxidation occurs every cycle; no special carbonation reset occurs during this block.

Measure cycle-resolved ethylene and cumulative output throughout, including a fixed terminal window \emph{before any special reset}. Report the steam-minus-Ar output difference and its uncertainty on equal initial mass and feed bases. This is operating performance, including any steam carryover or reversible inhibition, not proof of permanent damage. Do not insert an extra dry conditioning step before measuring it. Then apply the pilot-qualified terminal reset below once to each completed steam and Ar history and measure a common dry-cycle functional block. Compare the pre-reset and post-reset outputs with their matched Ar history: a steam penalty before treatment followed by recovery is distinct from tolerance during ordinary cycling. Recovery means performance under that terminal treatment, not proof of a unique repair mechanism. The pre-reset blocks also supply the untreated policy comparator for Experiment~2.

If a functional or material contrast warrants diagnosis, use the two five-minute segments in Table~\ref{tab:cyclic-timing} on separate specimens at 700\,$^\circ$C and atmospheric pressure. Ten mol\% steam and 5 mol\% added CO\textsubscript{2} are initial diagnostic settings, not an optimum or a simulation of nearly pure steam. Their results do not substitute for the ordinary-cycle comparison at the selected steam-rich duty. Keep total inlet molar flow, temperature history, pressure, specimen mass, and total elapsed exposure fixed. Calibrate switching holdup using an empty reactor and the actual water/CO\textsubscript{2} signals. Match the \emph{integrated inlet} CO\textsubscript{2} dose within each timing pair and report measured background CO\textsubscript{2}; equal setpoints alone are insufficient. Uptake and endogenous release remain outcomes rather than quantities forced to match.

\begin{table}[ht]
\centering
\caption{Timing and steam are crossed to distinguish a steam-dependent timing effect from ordinary dry-gas history. All unspecified gas is Ar.}
\label{tab:cyclic-timing}
\begin{tabularx}{\linewidth}{@{}lXX@{}}
\toprule
Arm & First five minutes & Second five minutes\\
\midrule
Wet, concurrent (P) & 10\% steam + 5\% CO\textsubscript{2} & Dry Ar\\
Wet, delayed (R) & 10\% steam & 5\% CO\textsubscript{2}, dry\\
Dry, early CO\textsubscript{2} & 5\% CO\textsubscript{2}, dry & Dry Ar\\
Dry, late CO\textsubscript{2} & Dry Ar & 5\% CO\textsubscript{2}, dry\\
Ar reference & Dry Ar & Dry Ar\\
\bottomrule
\end{tabularx}
\end{table}

The two dry timing arms are necessary because early CO\textsubscript{2} could protect against a dry exposure independently of steam. After the common reset described below, measure the mean ethylene amount per cycle and initial catalyst mass, $q$, over a fixed block for each arm. Define the steam-dependent timing contrast as
\begin{equation}
 D_{\mathrm{steam}}=(q_{\mathrm{wet,P}}-q_{\mathrm{wet,R}})
 -(q_{\mathrm{dry,early}}-q_{\mathrm{dry,late}}).
 \label{eq:cyclic-interaction}
\end{equation}
A positive wet P--R difference establishes a timing benefit under steam; positive $D_{\mathrm{steam}}$ adds evidence that steam changes that benefit. Neither alone proves hydrolysis or lithium volatilization. Also report $q_{\mathrm{wet,P}}-q_{\mathrm{dry,early}}$ and $q_{\mathrm{wet,R}}-q_{\mathrm{dry,late}}$, with uncertainty intervals against the predeclared consequential loss. These compare steam substitution within each CO\textsubscript{2} schedule. Report every arm's absolute yield and its difference from Ar as an overall history comparison: both wet arms might be damaged equally even when their timing interaction vanishes. A null here concerns function after the specified CO\textsubscript{2} history and reset, not untreated steam compatibility.

Use independently run sister specimens, randomized within preparation and run-order blocks, for the first-segment endpoint, completed exposure, and final reset endpoint. At every endpoint, use in situ observation or a validated termination, quench, and controlled transfer, documenting the elapsed time and gas history. The first-segment samples can locate changes already present during steam only for observables shown to survive sampling. Withhold temporal species or spatial attribution if termination changes the measured state. Final specimens alone cannot distinguish steam damage from changes during the later Ar/CO\textsubscript{2} segment. Repeated cycles on one specimen do not replace independent runs. Determine the replicate count from pilot variance and a predeclared minimum consequential yield or lithium-inventory difference.

\subsection{A common reset, defensible inventories, and a persistent function assay}

Use one common terminal carbonation/oxidation sequence for the recovery comparison and for all timing arms. It is additional to ordinary cycle oxidation and follows the ordinary comparison's pre-reset output measurement. In a pilot, choose dry-inlet CO\textsubscript{2}/O\textsubscript{2} conditions and a duration sufficient for the slowest responding specimen to reach stable net uptake and a reproducible carbonate-sensitive spectrum. Freeze that sequence for the comparison. A joint CO\textsubscript{2}/O\textsubscript{2} endpoint limits a new decarbonation interval while reoxidizing the carrier; its adequacy must be demonstrated on this formulation. Measure CO\textsubscript{2}, O\textsubscript{2}, H\textsubscript{2}O, CO, and H\textsubscript{2} throughout. Apply the same state-preservation qualification to phase and Fe-state measurements after the reset \emph{and after the actual washout}.

The objective is a demonstrably completed treatment of the retained carbonatable fraction, not equal absolute carbonate inventory if Li has been lost. An uptake plateau can reflect slow transport or passivation, so it needs the independent species/state check. If a shared process-compatible reset cannot establish interpretable endpoints, report the resulting recovery limit without asserting absolute irreversibility. Any further reset needed to assess incomplete recovery must be applied comparably and charged to the recovery policy.

For the post-reset assay, verify CO\textsubscript{2} and O\textsubscript{2} washout before introducing ethane against thresholds qualified to exclude a consequential change in the first ethane pulse. Gao reports CO\textsubscript{2} inhibition on coated LSF in Fig.~4D; the explicit reversible switch trace in Fig.~4E is a separate molten-Li\textsubscript{2}CO\textsubscript{3}/ethane/O\textsubscript{2} experiment.\cite{gao2020} Residual gaseous O\textsubscript{2} can also change output independently of carrier oxygen. Establish the maximum washout needed across arms and apply the same interval in the diagnostic; then determine whether that interval alters the coating contrast. An effect erased by washout cannot be claimed as persistent protection. Follow identical ethane/oxidation cycles with dry Ar purges for a fixed block selected in the pilot to distinguish gas carryover from retained function.

Measure ethylene directly as the primary outcome, separately from the source's C\textsubscript{2+} objective, with conversion, total CO\textsubscript{x}, other hydrocarbons, and retained carbon. Include carbonate-carbon capture, release, and inventory changes in the cycle carbon balance. Use these measurements to bound the carbonate contribution before assigning hydrocarbon-derived CO\textsubscript{x}; regeneration CO\textsubscript{2} alone is not a coke assay. If carbonate and coke contributions cannot be resolved, report total CO\textsubscript{x} and the attribution limit rather than a precise combustion selectivity or coke amount. The close regeneration precedent also identifies this accounting problem.\cite[ch.~4.1]{chacko2025thesis} No isotope campaign is required for the direct ethylene-output decision. Use a thermal blank to bound noncatalytic ethane conversion. Water and H\textsubscript{2} must be measured directly: carbonate/hydroxide and carrier oxygen changes make water alone an ambiguous lattice-oxygen reporter.

Close the Li inventory over active particles, grit, accessible reactor walls, collected particulates, and downstream condensates. Use complete digestion with blanks and recovery spikes for each matrix; establish sampling representativeness and extraction recovery before aging. Record each exposure segment's capture separately, but recognize that earlier deposits may be released later. Bulk elemental assays locate total Li; ordinary EDX cannot establish Li coverage.

If redistribution controls the interpretation, a candidate is cross-sectional NanoSIMS mapping of Li-rich patches and rim-to-interior gradients, correlated with electron microscopy. Qualify submicrometre spatial resolution finer than the features being compared, matrix-dependent ion yields, sectioning artifacts, and beam-induced redistribution. Map intensities do not by themselves give absolute coating coverage, and unresolved films remain outside the claim. For carrier state, use phase characterization and, subject to access and qualification, \textsuperscript{57}Fe M\"ossbauer spectra on state-preserved sister samples, following the method precedent in Gao.\cite{gao2020} Reference materials and replicate spectra must distinguish the proposed Fe-state or phase change above fitting uncertainty; overlapping components and bulk averaging can conceal an interfacial change. Carbonate-sensitive vibrational analysis supplies a separate chemical indicator. Instrument access and adequate resolution are dependencies, not established capabilities; close the corresponding mechanism claim if they cannot be qualified.

The analytical scale is demanding. On an explicitly assumed final-catalyst basis, 1\,g containing 10\,wt\% Li\textsubscript{2}CO\textsubscript{3} has about 18.8\,mg Li. A 0.1\% change is 18.8\,$\mu$g, below a 1\% bulk-inventory error. At 100 standard cm\textsuperscript{3}/min for five minutes, 5\% CO\textsubscript{2} supplies approximately 1.12\,mmol. Conversion of 0.1\% of the assumed carbonate inventory corresponds to only 1.35\,$\mu$mol CO\textsubscript{2}, about 61\,ppm of the total gas dose against a 50\,000\,ppm feed (0\,$^\circ$C standard-volume convention). These are resolution examples, not predicted losses. Captures, absolute assays, and chemical signatures must complement subtraction of large inlet/outlet signals.

If a single diagnostic exposure is unresolved, choose the number of cycles from the full-method uncertainty budget, then accumulate captures over that fixed block. Include actual ethane and ordinary reoxidation intervals. Apply the assigned two-segment timing history once per cycle but the special common reset only once, after the accumulation block; freeze and report this cadence before assigning a history-specific loss rate. This ten-minute diagnostic exposure remains distinct from the five-minute ordinary purge. Increasing temperature or extending steam exposure until damage appears would test a different condition, not establish failure of the proposed purge. A useful null excludes consequential change only under its stated gas, duty, cycle count, catalyst state, and recovery policy; it is not proof of lifetime steam tolerance.

\subsection{Experiment 2: test a usable policy and its complete-cycle value}

A consequential ordinary-cycle penalty, a recovery benefit, or a timing contrast justifies comparing management policies at the same selected steam-rich duty. Retain the post-ethane Ar purge. Freeze steam dose, CO\textsubscript{2} circulation limit, total flow, temperature trajectory, and treatment cadence before testing; any unavoidable change in steam delivery is part of the policy comparison. The diagnostic ratio cannot be transferred automatically: 10\% steam/5\% CO\textsubscript{2} gives a ratio of two; maintaining that ratio in a binary purge requires one-third CO\textsubscript{2}. This is a composition calculation, not a required protective dose. At the same five-minute total purge flow and one-minute 81.6\% ethane feed used in the source cycle, one-third CO\textsubscript{2} would circulate about 2.04 mol CO\textsubscript{2} per mol ethane fed in that one purge. Preserving steam dose instead of total flow increases gas throughput and requires a new heat and hydrodynamic comparison.

Compare concurrent protection, delayed restoration, untreated steam without supplementary CO\textsubscript{2}, and dry Ar over independently replicated operating blocks. The pre-reset ordinary-cycle comparison supplies the latter two policies when conditions, block length, and preparation/run blocks match; otherwise include matched controls. Untreated steam and Ar receive only ordinary oxidation during the policy block. Fix any periodic CO\textsubscript{2} treatment in the managed policies from the pilot and record its cadence explicitly; a terminal analytical reset occurs only after the counted block and does not improve its reported output. Delayed restoration must be allowed a demonstrated process-compatible treatment, rather than intentionally insufficient recovery. If actual Li export is consequential and retained-material restoration fails, test one measured replenishment operation as a competing policy. Replacing Li can also change coverage, salt composition, and carrier structure, so it is not automatically a selective intervention on Li inventory. Count its interruption, heating, washout, and achieved recovery; do not infer an online replacement method from the original synthesis.

Let $q_i$ be block-averaged desired product per cycle, $\tau$ the base cycle time, and $h_i$ the average added treatment time, including periodic operations. For equal initial catalyst mass and feed pulses,
\begin{equation}
 J_i=\frac{q_i}{\tau+h_i},\qquad
 J_P>J_R\ \Longleftrightarrow\
 h_P<\frac{q_P}{q_R}(\tau+h_R)-\tau.
 \label{eq:cyclic-value}
\end{equation}
Use integrated measured product over the entire block, including aging, rather than a recovered peak. As an illustration only, fully restoring the published Ar-aged C\textsubscript{2+} yield from 47.24\% to 53.2\% permits at most
\begin{equation}
 16\left(\frac{0.532}{0.4724}-1\right)=2.02\ \text{additional minutes per cycle}
 \label{eq:cyclic-threshold}
\end{equation}
before losing the throughput advantage, even before charging gas and heat. A five-minute added step every cycle needs a 31.25\% output increase; the illustrated complete recovery provides 12.62\%. This is a threshold for the documented Ar loss, not a prediction of steam damage or a full-life prevention benefit.\cite{brody2022}

Report ethylene and C\textsubscript{2+} productivity separately per initial catalyst inventory and elapsed time, together with net fresh and circulating CO\textsubscript{2}/steam, Li handling, and heat requirements per product. Recycling can reduce fresh CO\textsubscript{2} use while leaving circulation and separation burdens. A gas-intensive protection policy may remain a tradeoff, but a smaller Li signal alone cannot establish process value.

\subsection{Experiment 3: predict recoverability under a new operating history}

The scientific extension is to identify a measured state indicator that predicts persistent function under the specified recovery policy. If Li export is a leading candidate, calibrate per-cycle Li transfer for the protected and delayed policies and relate retained inventory to subsequent ethylene output on independently aged specimens. Test whether that relationship is single-valued within uncertainty. If equal retained Li gives different function, an inventory-only model fails; assess an independently measured coating-distribution or carrier-phase indicator rather than fitting a separate loss constant for each history. Successful prediction establishes usefulness over the tested range, not causal dominance of Li inventory over correlated coating or carrier changes.

Use continuous protected and continuous delayed operation as calibration histories, then withhold one alternating protection/recovery schedule for validation. Fix its period, operating-treatment cadence, prediction model, and uncertainty intervals before collecting it. Predict cumulative Li transfer and cycle-resolved useful output through the operating block, plus standardized function after its single terminal diagnostic reset. Do not insert that terminal reset between cycles unless it is explicitly part of the operating policy and charged there. An additive per-cycle loss model supplies a clear initial prediction; failure under alternation reveals history dependence that the simple inventory description misses. Add a measured state only if it explains and predicts that failure on a further independent block. The result must improve upon a carrier-state-only model and a simple cycle-count aging description sufficiently to change a recovery or purge decision. A good fit to the calibration histories is not that result.

This sequence can yield a substantial program if the selected steam duty changes useful function and a measured state predicts recovery. An ordinary-cycle steam-minus-Ar interval excluding consequential loss supports compatibility without supplementary CO\textsubscript{2} only for that composition, duty, cycle count, conditioning, and ordinary oxidation policy; it closes the need to claim protection against an undetected penalty in that test. A pre-reset loss followed by terminal recovery instead supports recoverability under that treatment and leaves its operating cost to Experiment~2. Close the claim of persistent concurrent-over-delayed benefit after the stated reset only when the uncertainty interval for the predeclared direct contrast $q_{\mathrm{wet,P}}-q_{\mathrm{wet,R}}$ excludes a consequential advantage. Use the interval for $D_{\mathrm{steam}}$ to assess whether steam modifies that timing benefit. Individual recovery bounds remain separate conclusions, and neither comparison establishes that CO\textsubscript{2} treatment is unnecessary. If protection preserves function but cannot repay its complete-cycle burden against untreated steam and Ar, retain the scientific limit and reject the practical intervention. No result here establishes post-ethane steam compatibility or lifetime tolerance. Explaining the original Ar-aged decline requires representative Ar-aged material and a separate recovery study; it must not be substituted silently for failure of the steam hypothesis.
